\section{Introduction}
\label{sec:intro}

The explosive growth of biomedical literature---over 35 million citations in PubMed alone---has created text repositories that no human can fully assimilate. Modern AI methods promise to turn this deluge into actionable scientific insight, yet two persistent gaps separate current capabilities from the goal of \emph{automated, trustworthy hypothesis generation}~\cite{zhang2025exploring,bhasuran2023literature}. First, large language models (LLMs) generate fluent but ungrounded claims: their outputs cannot be traced back to specific evidence paths, a failure mode known as hallucination~\cite{zhao2026survey,brown2025systematic}. Second, the symbolic knowledge graphs (KGs) that could anchor LLM reasoning---such as SemMedDB, which distills 33 million predications from PubMed~\cite{kilicoglu2012semmeddb,kilicoglu2020broadcoverage}---are themselves noisy, with SemRep precision ranging from 0.55 to 0.80 depending on predicate type. Naively feeding a noisy KG into an LLM amplifies, rather than mitigates, ungrounded generation.

Recent work on KG-augmented LLMs has made encouraging progress. Think-on-Graph (ToG)~\cite{sun2023thinkongraph} and its successor ToG~2.0~\cite{ma2024thinkongraph} use the LLM to actively explore KG neighborhoods, while Graph-constrained Reasoning (RoG)~\cite{luo2024graphconstrained} constrains LLM decoding with KG-grounded paths and GNN-RAG~\cite{mavromatis2025gnnrag} retrieves reasoning paths via graph neural networks. These methods explore KG paths during reasoning, but they treat the KG as a fixed, trusted resource: they do not \emph{refine} the KG itself, and they do not post-hoc verify whether the final output is entailed by a specific path. Meanwhile, the literature-based discovery (LBD) community has long sought to re-discover hidden disease--treatment links via co-occurrence and graph traversal~\cite{swanson1986undiscovered,hristovski2015constructing,bhasuran2023literature}, but classical LBD methods lack the generative flexibility of LLMs and produce rankings rather than explainable hypotheses. The gap between \emph{generative LLM reasoning} and \emph{verifiable symbolic grounding} thus remains open.

We argue that closing this gap requires \emph{LLM-based KG verification}: the LLM should not only consume the KG but also help clean it, and the KG should not only scaffold generation but also verify the output. Concretely, an LLM can verify each triple in a retrieved subgraph, retaining supported triples and deleting contradicted ones. This verified KG becomes the symbolic scaffold for generation, and the same paths that scaffold a hypothesis also serve as its traceable evidence.

The choice of LLM-based verification over alternative KG-cleaning approaches is deliberate. Traditional statistical outlier detection methods~\cite{yao2021typing} operate on graph structure alone, missing semantic plausibility: a triple like \texttt{(aspirin, TREATS, diabetes)} is structurally valid but medically incorrect, and only a model with biomedical knowledge can flag it. KG embedding models~\cite{trouillon2016complex,sun2019rotate} can score triple plausibility but, as we show empirically, their scores are unreliable when the training KG itself is noisy (MRR$<$0.05). LLMs, by contrast, bring both semantic understanding and broad biomedical knowledge to verification: they can assess whether a predicate is medically plausible given the entity types, detect direction errors (e.g., \texttt{(diabetes, TREATS, metformin)} vs.\ the correct \texttt{(metformin, TREATS, diabetes)}), and identify type violations. The cost is LLM API calls, which we bound by verifying only the retrieved subgraph ($K{=}40$ triples per query).

We introduce \textbf{KG-SCoRE} (\emph{Knowledge-Graph-Symbolic Co-Refinement}), a pipeline that realizes this loop. Module~A verifies KG triples via LLM prompting (optionally gated by ComplEx~\cite{trouillon2016complex} scores). Module~B retrieves multi-hop subgraphs. Module~C conditions LLM hypothesis generation on the subgraph, forcing every hypothesis to carry an explicit KG path. Module~D verifies each hypothesis's entailment against its path, reporting both strict and broad faithfulness. While RoG constrains generation to KG-grounded paths during decoding, KG-SCoRE adds post-hoc verification of whether the final output is entailed by a specific path---a dimension not reported by prior methods.

We evaluate KG-SCoRE on a 9{,}188-triple biomedical KG with 20 LBD gold pairs. Under fair evaluation (all methods receive post-hoc path extraction), KG-SCoRE achieves StrictFaith@10 of 0.10 (vs.\ ToG's 0.05 and RoG's 0.04) and BroadFaith@10 of 0.99. On the 10-pair multi-seed subset ($n{=}30$), no significant difference from RoG is detected (McNemar $p{>}0.20$), while KG-SCoRE significantly outperforms ToG ($p{<}0.05$), Vanilla LLM ($p{<}0.01$), and BM25-RAG ($p{<}0.05$). KG cleaning evaluation shows the LLM verifier achieves perfect noise recall (1.0) but low precision (0.51), explaining why aggressive cleaning shifts the signal-to-noise ratio favorably. Ablations reveal that ComplEx-based learned scoring \emph{harms} recall on noisy KGs (MRR=0.013), confirmed by $\tau$ sensitivity analysis across four thresholds, while LLM-only verification achieves both competitive recall and near-perfect traceability.

Our contributions are:
\begin{itemize}
  \item An \emph{LLM-based KG verification} approach (Module~A) that cleans noisy triples by prompting the LLM to verify each retrieved triple, with direct cleaning accuracy evaluation.
  \item A \emph{symbolic-scaffolded generation} paradigm (Module~C) with a \emph{faithfulness@\textit{k}} metric (Module~D) split into strict and broad variants, evaluated fairly across all methods via post-hoc path extraction.
  \item An empirical evaluation revealing that LLM verification achieves competitive recall \emph{and} higher strict faithfulness than baselines, while learned KG-completion scoring is counterproductive on noisy KGs---a design insight for KG-augmented LLM systems.
\end{itemize}
